By \(\texttt{def:pair-action}\), the nonidentity element \(s_b\) interchanges the two members of block \(b\), preserves the arm label, and fixes every member outside that block. Hence, for \(r,a\in\{0,1\}\),
\[
s_b(v_{b,r},a)=(v_{b,1-r},a),
\]
and the two arm-specific orbits in block \(b\) are exactly
\[
O_{b,0}=\{(v_{b,0},0),(v_{b,1},0)\},\qquad O_{b,1}=\{(v_{b,0},1),(v_{b,1},1)\}.
\]
The action on complete records is the relabelling specified by \(\texttt{def:group-action}\).

Suppose the target arm-\(1\) lifted law is supported on \(O_{b,1}\). Its orbit-mass vector assigns all its mass to the arm-\(1\) orbit \(O_{b,1}\). An exclusive mate-control analogue instead takes values in the distinct arm-\(0\) orbit \(O_{b,0}\), so its corresponding orbit-mass vector assigns zero mass to \(O_{b,1}\). The vectors differ, and therefore the orbit equalities in \(\texttt{def:joint-orbit-balance}\), equivalently the relevant rows of \(Aw=b_{\mathrm{LP}}\) in \(\texttt{def:common-selector-feasibility-lp}\), cannot hold for that analogue. There is also a direct legality failure: \(\texttt{ass:observable-selector-kernel}\) requires an arm-\(1\) selector to be supported on members whose realized arm is \(1\), and thus on \(\mathcal U_1\), whereas a mate-control state lies in \(\mathcal U_0\).

Now consider separately a legal implemented arm-\(1\) selector with \(q^{\mathrm{obs}}_{w,1}(O_{b,1})=0\). Because the target arm-\(1\) law is supported on \(O_{b,1}\), every target-supported arm-\(1\) orbit has zero calibration mass. This is exactly the premise of \(\texttt{prop:disjoint-orbit-witness}\) with \(a=1\). If \(k\le M\), that proposition supplies its zero-coverage construction. If \(k=M+1\), \(\texttt{prop:rank-boundary}\) instead gives \(I_{0,w}=I_{1,w}=[-K,K]\), \(\mathcal C_w=[-2K,2K]\), and coverage one. These two cases exhaust the possible rank regimes.